% Autogenerated translation of COVER_LETTER_IJIS.md by Texpad
% To stop this file being overwritten during the typeset process, please move or remove this header

\documentclass[12pt]{book}
\usepackage{graphicx}
\usepackage{fontspec}
\usepackage[utf8]{inputenc}
\usepackage[a4paper,left=.5in,right=.5in,top=.3in,bottom=0.3in]{geometry}
\setlength\parindent{0pt}
\setlength{\parskip}{\baselineskip}
\setmainfont{Helvetica Neue}
\usepackage{hyperref}
\pagestyle{plain}
\begin{document}

\chapter*{Cover letter — International Journal of Information Security}

Dear Editor-in-Chief,

We submit the manuscript "Glaukopis: Auditable Graph-Grounded Instruction Tuning for Cyber Threat Intelligence" for consideration as a regular article in the International Journal of Information Security.

\textbf{What the paper contributes.} Security operations centers that cannot send incident data to hosted models need to specialize an open-weight language model for cyber threat intelligence (CTI) tasks, and they need to know where the training data came from and what the specialization cost. The paper presents (i) a data-construction method in which every CTI instruction-tuning target is the deterministic rendering of a typed query over a knowledge graph of twelve public CTI sources, so that each training example resolves to the public facts that justify it; (ii) a retained-checkpoint evaluation protocol with item-level bootstrap intervals and response-level error taxonomies that locates where a 32B model gains and loses capability across a four-stage fine-tuning recipe; and (iii) a quantified audit of training–evaluation overlap at the n-gram, entity and gold-fact levels. The evaluation covers three CTI benchmarks, a general-reasoning diagnostic, the base model, a security-specialized open model, six hosted frontier models and a single-stage fine-tuning run. The results include substantial negative findings — a 17,320-row structured-answer stage collapses actor attribution and biases CVSS exploitability judgments, and general reasoning falls by 11.6 points — and the paper explains both mechanistically from the content of the training shards. We believe this combination of a traceable data method and an honest, diagnostic evaluation fits the journal's scope and its readership of security practitioners and researchers.

\textbf{Prior history.} An abridged conference version was submitted to IEEE ICTAI 2026 and the full version to \emph{Computers \& Security} (COSE-D-26-05507); neither was accepted or published, and no preprint has been posted. The present manuscript addresses every point raised in both sets of reviews: it rewrites the exposition to be self-contained, states the rationale for each design choice, adds a baselines section and a single-stage reference run, adds mechanistic analyses of the two regressions, locates the general-reasoning loss by subject, reframes the cost analysis on durable measures, and supplies full citations for all reference models. We would be glad to provide the earlier reviews and our point-by-point response on request.

\textbf{Declarations.} The work was funded by Athena Security Group, which develops the system described; the corresponding author is a co-author of one of the three benchmarks used, and this developer–evaluator overlap is disclosed and analyzed in the manuscript (Section 7.3). No human participants or personal data were involved. Large language models were used only as writing and coding aids, as stated in Section 7.1 and the Declarations. The manuscript is not under consideration elsewhere, and all authors have approved it.

\textbf{Suggested reviewers.} [Optional: three or four independent researchers in LLMs for security or CTI benchmarking, with institutional email addresses; none affiliated with the authors, Athena Security Group, Florida Atlantic University EECS, or the AthenaBench team.]

Thank you for considering our work.

Sincerely,

Peter J. Worth Jr. (corresponding author)
Department of Electrical Engineering and Computer Science, Florida Atlantic University, Boca Raton, FL, USA
pworthjr@athenasecuritygrp.com
on behalf of Ionut Cardei

\end{document}
